\section{Related Work}\label{s:related}
Araújo et al. \cite{araujo_energy_2024} give an overview of research concerning energy consumption by microservice applications through an extensive literature review. While it was the most evaluation method used by most papers, the authors note that 10\% of papers that simulate microservices do not specify what simulator they use or how the microservices are simulated. In the papers that do disclose the simulators used, \textit{CloudSim} is by far the most frequently used, followed by the \textit{DRACeo} simulator, the \textit{Linpack} benchmarker and finally \textit{Matlab}. This study shows that simulators, particularly \textit{CloudSim} are utilized by researchers to investigate energy consumption of microservice applications. Furthermore, Araújo \etal note that authors seem to find the simulators used inconsequential to their research. The research discussed here builds on these observations by investigating what cloud simulators are well-suited to simulate microservice applications and showing how accurate they are. This also highlights that what simulator is chosen can indeed have an impact on the results obtained through its use.

Mansouri et al. \cite{mansouri_cloud_2020} present an extensive review of 32 cloud simulators including \textit{CloudSim}, \textit{NetworkCloudSim} and \textit{DISSECT-CF} describing their implementations, strenghts, limitations, and available features, including whether they provide energy modelling. They note the importance of including energy consumption simulation in cloud simulators and that the simulators that do provide this only do so in a quite simplified manner, often only taking CPU utilization into account. While they do consider metrics of the simulators themselves such as their time and energy consumption, they do not consider the accuracy of the simulated results. They do note that the simulators are not very scalable, which is an important capability with cloud systems rapidly increasing in scale. The research presented here builds on this paper by focusing on the energy consumption simulation capabilities of a smaller selection of cloud simulators.

Papers by Ismael \cite{ismail_energy-driven_2020} and  Vincenzo Emanuele Martone et al. \cite{martone_evaluation_2023} both describe cloud simulators with a focus on their abilities to simulate energy consumption metrics. They are both literature reviews that extract their information from other reviews, simulator papers or general properties of the simulators and do not test the simulators themselves as opposed to the simulated metrics as is done in the research discussed here. In \cite{ismail_energy-driven_2020}, eleven simulators including \textit{CloudSim} and \textit{DISSECT-CF} are discussed. Simulators are evaluated in terms of the hardware components for which energy consumption can be simulated, what energy metrics can be extracted, and energy management techniques that are provided to lower the energy consumption. Furthermore, the impact or popularity of the simulators is discussed, specifically whether the simulators were used or extended in subsequent works. Finally, Ismael \cite{ismail_energy-driven_2020} notes that most simulators only support energy simulation of certain components, not of complete systems through the entire energy chain. Furthermore, there is a lack of support for multi-objective simulation and optimization, as simulators that support energy simulation often do not support e.g. modelling of network communication.

Martone et al. \cite{martone_evaluation_2023} describe five simulators, again including \textit{CloudSim} and \textit{DISSECT-CF} according to their implementation and energy simulating capabilities and evaluates them according to when they were last updated, popularity, availability, granularity and performance profile. For the last metric, values such as simulators' run times and memory usage are taken from other papers. Finally, they recommend \textit{CloudSim} as the generally most suitable simulator for general-purpose simulation of cloud systems given its widespread adoption, extensive documentation and numerous available available extensions.

In a recent paper, Hasanein Rjeib and Gabor Kecskemeti \cite{rjeib_towards_2025} do perform evaluations of the accuracy of simulated energy consumption measures by \textit{CloudSim} and \textit{DISSECT-CF} using the PlanetLab workload traces. For these traces, they find that \textit{DISSECT-CF} simulates energy consumption more accurately, in part because it also provides the capabilities to simulate the energy usage of resources such as RAM. Furthermore, they observe some abnormal behaviour by CloudSim when CPU usage is 0 for a VM. They note that, while CloudSim is simpler to implement and easier to use, DISSECT-CF has a higher accuracy and lower simulation time, making it overall the preferred simulator. This paper is somewhat similar to the research described here but differs as it it only considers two simulators and does not test microservice applications. Moreover, it only very briefly discusses the usability of the simulators.

While previous work such as that discussed in this section has compared cloud simulators on a higher level, both for general use cases and focused on simulating energy usage, this study goes further by actually carrying out simulations and providing the produced energy values to real-world measures. Furthermore, this is done in the context of microservice applications, which is a context in which cloud simulators are used to support research, but are very rarely the subject of study themselves.

% \TODO{
% Describe here scientific papers similar to your experiment, both in terms of goal and methodology.  One paragraph for each paper (we expect about 5-8 papers to be discussed). Each paragraph contains: (i) a brief description of the related paper and (ii) a black-on-white description about how your work differs from the related paper. You may place this section immediately after the Background section, if necessary.
% }